\section{Fixed all-test strict sources need not have a density neighborhood of solvability}
\label{sec:r17-fixed-strict-source-nonopenness}

\paragraph{Status and precise purpose.}
This section proves a new original-category obstruction at a
\emph{fixed} reference source. It is not a new parameterized existence
theorem. The reference potential already permits both signs of every
original compact smooth scalar direction. Nevertheless arbitrarily
small strictly positive original smooth density changes can leave
the source image of the entire prescribed strong completion.
The proof uses remaining degree one, where a linear moment identity
is valid. No nonlinear Gauss-trace pushforward is asserted.

The new statements have status \texttt{NEW\_PROVED} relative to the
frozen R12--R16 chain, subject to the author-separated audit recorded
in the companion ledger. The existing simultaneous probe, canonical
convex Gauss traces, global Green energy, and literal fixed-source
domination are \texttt{INTERNAL\_INPUT}; they are not reproved as
new. The full completed moment-Laplacian identity itself is already
SO6, following DO6--DO9; the additional $H_0^1$ boundary and
contraction statement below is new. Ordinary real Green theory
and almost-everywhere differentiability of a monotone function
are explicitly identified \texttt{EXTERNAL\_KNOWN} inputs.

\subsection{One allowed original background, with all original tests retained}

Fix the prescribed finite extension $K/\mathbf Q_p$, its unchanged
absolute value, and an original uniformizer $\pi$. Set
$\lambda=-\log|\pi|>0$ and
\[
 \begin{split}
 X&=(\mathbf G_{m,K}^{2})^{\rm an},\qquad
 t=(\log|T_1|,\log|T_2|),\\
 Q&=(-\lambda,\lambda)^2,\qquad V=t^{-1}(Q),\\
 \beta&=d't_1\wedge d''t_1+d't_2\wedge d''t_2,
 \qquad S=\beta\wedge[X].
 \end{split}
 \tag{FP1}
\]
The closed box lies in the ordinary interior of the original compact
affinoid $t^{-1}([-2\lambda,2\lambda]^2)$, and $S$ extends to a
neighborhood of that larger box. All radii are in the original value
group. This is an allowed symmetric middle-positive separately
$d'$- and $d''$-closed background, here in fact supplied by a smooth
strongly positive form. This additional description is specific to
the counterexample; no general positivity cones are identified.
Here $n=2$, $p=1$, $q=1$, and $d=2$.

Let $\mathscr C$ be the unchanged relative Borel strong completion
with its printed approximation boundary condition, and let
$\mathbb D=\mathscr C-\mathscr C$. Its quadratic mixed energy is
$B$, so $I(a,b)=B(a-b,a-b)$. Write $\mu_u$ for its completed
source measure. Actual defining functions retain boundedness where
used, finite mass, compact strict negative sublevels, and the original
canonical no-loss conditions. No ordinary boundary value is simply
postulated for an arbitrary completed member.

Let $\sigma:Q\to V$ be the original joint Gauss section. Frozen
NF supplies a canonical finite nonpositive convex trace $g_u$ for
each $u\in\mathscr C$. It is additive, homogeneous, locally uniform
under strong convergence, and agrees with $u\circ\sigma$ only
Lebesgue-almost everywhere for a general relative representative.
All arguments below use this canonical trace; no point value of a
Borel class at an outer-null Gauss point is chosen.
The printed NF unit box is replaced by $Q$ using the same proof
after a real affine coordinate rescaling; the original field,
absolute value and moment map are unchanged. Formula (FP6) below
independently fixes the energy normalization in the present coordinates.

\subsection{The degree-one moment equation and its completed zero boundary}

\begin{lemma}[A continuous zero-boundary distribution with an energy bound]
Let $Q\subset\mathbf R^2$ be a bounded open rectangle. Suppose
$g\in C(\overline Q)$, $g|_{\partial Q}=0$, and for some $C<\infty$,
\[
 \left|\int_Qg\,\Delta\varphi\,dx\right|
                   \le C\|\nabla\varphi\|_{L^2(Q)}
                    \quad(\varphi\in C_c^\infty(Q)).
 \tag{FP2}
\]
Then $g\in H_0^1(Q)$ and $\|\nabla g\|_2\le C$.
\end{lemma}

\begin{proof}
Extend $g$ continuously by zero outside $Q$. Let $Q_\varepsilon$
be inner rectangles exhausting $Q$, and mollify with a nonnegative
even kernel of radius less than half the distance of
$Q_\varepsilon$ to $\partial Q$, obtaining $g_\varepsilon$.
For a smooth test on $Q_\varepsilon$, move the convolution onto
the test in (FP2). Its support stays in $Q$, and convolution
does not increase its $L^2$ gradient norm. Density therefore gives
\[
 \|\Delta g_\varepsilon\|_{H^{-1}(Q_\varepsilon)}\le C,
 \tag{FP3}
\]
where $H_0^1$ is normed by the gradient.

Let $H_\varepsilon$ be the harmonic extension of the smooth
boundary values $g_\varepsilon|_{\partial Q_\varepsilon}$.
Rectangles are regular for the ordinary Dirichlet problem. Equivalently
one may minimize the real Dirichlet energy with those smooth trace
values; the weak maximum principle and the elementary rectangle
barriers give its continuous boundary values. Then
$v_\varepsilon=g_\varepsilon-H_\varepsilon\in
H_0^1(Q_\varepsilon)$ and
$\Delta v_\varepsilon=\Delta g_\varepsilon$. Testing this equality
by $v_\varepsilon$ in (FP3) gives
$\|\nabla v_\varepsilon\|_2\le C$.
Extend $v_\varepsilon$ by zero to $Q$ and extract a weakly
convergent subsequence in $H_0^1(Q)$.

Uniform continuity of the zero extension of $g$ implies
$\sup_{\partial Q_\varepsilon}|g_\varepsilon|\to0$.
The maximum principle gives the same bound for
$\|H_\varepsilon\|_\infty$. Also $g_\varepsilon\to g$ uniformly
on compact subsets. Hence $v_\varepsilon\to g$ locally uniformly,
identifying its weak $H_0^1$ limit with $g$. Weak lower
semicontinuity proves the claimed energy bound. This argument
is the step excluding an uncontrolled harmonic boundary summand.
\end{proof}

\begin{proposition}[The completed moment identity retains the real zero boundary]
In (FP1), every $u\in\mathscr C$ satisfies
\[
 g_u\in H_0^1(Q),\qquad
 \|\nabla(g_u-g_v)\|_2^2\le I(u,v)
                              \quad(u,v\in\mathscr C),
 \tag{FP4}
\]
and
\[
                         t_*\mu_u=\Delta g_u
                   \quad\text{as Radon measures on }Q.
 \tag{FP5}
\]
Formula (FP5) is the frozen SO6 identity in this normalization.
The new conclusion is (FP4), including its boundary condition.
In particular, equality of two prescribed original source measures
fixes the corresponding real zero-boundary Poisson equation uniquely.
\end{proposition}

\begin{proof}
For an actual bounded admissible potential $a$ and
$\varphi\in C_c^\infty(Q)$, original current recursion, compact
Stokes and the full original top-form Gauss calibration give
\[
 \begin{split}
 \int_V\varphi(t)\,d\mu_a
  &=\langle aS,D(\varphi\circ t)\rangle
    =\int_Qg_a\Delta\varphi\,dx,\\
 B(a,\varphi\circ t)&=-\int_Qg_a\Delta\varphi\,dx,\qquad
 B(\varphi\circ t,\varphi\circ t)
                         =\int_Q|\nabla\varphi|^2dx.
 \end{split}
 \tag{FP6}
\]
The factor is $(n-1)!=1$ with the unnormalized $\beta$ in (FP1).
The test $\varphi\circ t$ has compact original support because
$t$ is proper over compact subsets of $Q$. It is an exact difference
of admissible actual potentials by frozen HC14--HC15. Consequently
the full difference energy and its Cauchy--Schwarz inequality apply
to it. These identities use actual original tests and the fixed
smooth background, and require no unresolved delta-action quotient.

For two bounded actual admissible potentials $a,b$, put
$h=a-b$ and $g=g_a-g_b$. Each individual convex trace is continuous
in $Q$ and tends uniformly to zero at its boundary: otherwise a
sequence with value at most $-\epsilon$ would lie in the fixed
compact strict negative sublevel of the original potential while
its moments leave every compact subset of $Q$. Thus $g$ extends
continuously by zero to $\partial Q$. Formula (FP6) and the
positive energy matrix give (FP2) with $C=B(h,h)^{1/2}$.
The preceding lemma proves
\[
                 \|\nabla(g_a-g_b)\|_2^2\le B(a-b,a-b).
 \tag{FP7}
\]
In particular each actual trace lies in $H_0^1(Q)$.

Every completed member has a bounded actual strongly defining
row by the frozen clipping diagonal. In degree one, (FP7) makes
its traces Cauchy in $H_0^1(Q)$. Their NF locally uniform limit
is precisely $g_u$, so it is also the $H_0^1$ limit. Applying
this to two rows proves (FP4) with the original completed boundary
condition, rather than imposing a new boundary representative.
Pass in the first line of (FP6) using completed source continuity
on the fixed original compact test and the local convergence of
the traces. This proves (FP5). No expression involving a higher
power of $Dg_u$ has been used.
\end{proof}

\subsection{A real Poisson perturbation that destroys convexity}

We record the real lemma with all hypotheses needed for a possibly
nonsmooth reference trace. The reference Hessian may have a singular
positive part; no smoothness of the reference source density is assumed.

\begin{lemma}[Small smooth density changes of a locally strict convex source]
Let $g\in H_0^1(Q)$ be finite and convex on a rectangle $Q\subset
\mathbf R^2$. Put $\tau=\Delta g\ge0$. Suppose on a nonempty ball
$B\Subset Q$ that
\[
                         \tau\ge c\,dx\quad(c>0).
 \tag{FP8}
\]
There are $k_j\in C_c^\infty(B)$, $0\le k_j\le1$, and positive
$\epsilon_j\downarrow0$ after a subsequence, such that the unique
$H_0^1(Q)$ solution of
\[
                 \Delta v_j=(1+\epsilon_j k_j)\tau
 \tag{FP9}
\]
is not convex. Each $1+\epsilon_jk_j$ lies in $[1,3/2]$ after
discarding finitely many terms.
\end{lemma}

\begin{proof}
Fix a unit direction $e$ and an interior line segment in $B$
parallel to $e$. The convex restriction of $g$ to this segment
has a finite second-order expansion at almost every point:
its right derivative is monotone and is differentiable with finite
derivative almost everywhere; integrating that first-order
expansion gives the second-order expansion of the restriction.
Choose such a point $x_0$, so that explicitly
\[
 g(x_0+se)=g(x_0)+a s+\tfrac12 b s^2+o(s^2),
                         \qquad b\ge0.
\]
Put $A=b+1>0$. Thus $b$ is the second derivative, not the
coefficient of $s^2$. This one-dimensional fact
suffices; no second-order expansion in every direction is needed.
Let $G_Q(x,y)$ be the positive Dirichlet Green function for
$-\Delta$ on $Q$, and put
$K(y)=\partial_{ee,x}G_Q(x_0,y)$ for $y\ne x_0$.
With polar coordinates around $x_0$ and angle measured from $e$,
the fundamental solution and its harmonic correction give
\[
 K(y)=\frac{\cos(2\theta)}{2\pi r^2}+O(1),\qquad
                    r=|y-x_0|\downarrow0.
 \tag{FP10}
\]
Here the correction is uniformly bounded after two $x$ derivatives
on one fixed inner ball: its boundary data stay uniformly bounded
as $y$ ranges in that inner ball, and the maximum principle and
interior harmonic derivative estimate apply. In particular this
statement does not require a smooth boundary at the corners of $Q$.

Choose $R>0$ with $\overline B(x_0,R)\subset B$, small enough
that the leading term in (FP10) dominates the bounded remainder
on the fixed angular sectors $\cos(2\theta)\ge1/2$.
For $\delta_j\downarrow0$ take smooth radial cutoffs $\chi_j$
supported in $\delta_j<r<R$, equal to one on
$2\delta_j\le r\le R/2$, and between zero and one.
Choose a fixed smooth angular function $\zeta$ on the unit circle,
with $0\le\zeta\le1$, supported where $\cos(2\theta)>1/2$,
and equal to one on a smaller nonempty angular arc. Define
\[
 k_j(y)=\chi_j(r)\zeta((y-x_0)/r),
 \qquad k_j=0\text{ off that annulus}.
 \tag{FP11}
\]
These are smooth compact functions, with $0\le k_j\le1$ and
$Kk_j\ge0$ pointwise. On the smaller arc and middle annulus,
(FP10) gives $Kk_j\ge c_1r^{-2}$ for a fixed $c_1>0$.
Thus (FP8) implies
\[
 M_j:=\int_QK(y)k_j(y)\,d\tau(y)
       \ge c_2\log\frac{R}{4\delta_j}\longrightarrow+\infty.
 \tag{FP12}
\]
Any singular positive part of $\tau$ contributes with the same
nonnegative sign. For each fixed $j$, this integral is finite
because its support is compact and separated from $x_0$.

The distribution $k_j\tau$ belongs to $H^{-1}(Q)$: multiplication
of an $H_0^1$ test by the fixed compact smooth $k_j$ is bounded in
$H_0^1$, and $\tau=\Delta g\in H^{-1}$. Let $z_j\in H_0^1(Q)$
be its unique zero-boundary Poisson potential. It is harmonic near
$x_0$, and the Green formula, differentiated on that source-free
ball, gives
\[
                  \partial_{ee}z_j(x_0)=-M_j.
 \tag{FP13}
\]
This formula also holds for the present finite signed measure
on its compact support: convolve it there in $H^{-1}$, use the
smooth-source Green formula, and pass to the limit. The solutions
converge in $H_0^1$; on the source-free ball harmonic interior
estimates give convergence of every derivative, while the smooth
kernel on the compact source support passes through the measure
limit. This justifies (FP13) without pairing a singular kernel
with its pole.

Set $\epsilon_j=A/M_j$ and, if desired, pass to a subsequence
making these numbers decrease. Then $\epsilon_j\to0$ and the
unique solution of (FP9) is $v_j=g+\epsilon_jz_j$. Since
$k_j\tau\ge0$, the weak maximum principle also gives
$z_j\le0$, so this candidate remains nonpositive and retains
the zero boundary; the failed condition is convexity.
Its restriction to the chosen line has a second-order expansion
at $x_0$ with second coefficient $b-A=-1$. A convex function
cannot have such a negative coefficient, by its centered second
difference along $e$. Hence $v_j$ is not convex. Positivity of the density follows
for sufficiently large $j$ from $\epsilon_j<1/2$.
\end{proof}

\subsection{The fixed original source counterexample}

Say that a bounded actual admissible $w<0$ has the full
two-sided compact test property if
\[
 \text{for every real original compact smooth }\phi
 \text{ there is }s_\phi>0\text{ with }
       w+s\phi\text{ actual admissible whenever }|s|\le s_\phi.
 \tag{FP14}
\]
All actual refinements of $\phi$ are included. This condition is
nonempty: the frozen NP probe $\psi_0$ has it, is continuous and
bounded between $-1$ and zero, has finite mass and all finite
positive power energies, and retains its strict negative-sublevel
and no-loss boundary conditions. It is not a Banach interior claim.

\begin{theorem}[Failure of local density solvability at every such fixed source]
In the one fixed original geometry (FP1), let $w$ be any bounded
actual finite-mass reference with (FP14) and the printed actual
boundary and no-loss hypotheses. Then the following statements hold.
\begin{enumerate}
\item For every nonempty real ball $B\Subset Q$, there are
strictly positive original smooth densities
$f_j=1+\epsilon_j k_j\circ t$, with $k_j\in C_c^\infty(B)$,
such that
\[
  1\le f_j\le\tfrac32,\qquad
  \|f_j-1\|_\infty\longrightarrow0,
  \qquad f_j\mu_w\notin\{\mu_u:u\in\mathscr C\}.
 \tag{FP15}
\]
In particular the reference $w$ is fixed throughout the sequence.
\item Every $f_j\mu_w$ is a positive original Radon datum,
literal and strongly continuous on the entire completed difference
quotient. The data share the same null ideal as $\mu_w$, and
\[
 \|\ell_{f_j\mu_w}-J(w)\|_*
                   \le\|f_j-1\|_\infty\|J(w)\|_*
                       \longrightarrow0.
 \tag{FP16}
\]
They have one common full original tester capacity bound and one
common fixed-source strong $L^1$ modulus.
\item There is also a sequence as in (FP15) which tends to one
in the original smooth compact-open topology on $V$, as well as
uniformly. For this additional assertion the compact supports of
$f_j-1$ escape every compact subset of $V$; no common compact
support is asserted.
\end{enumerate}
\end{theorem}

\begin{proof}
First observe why the real lemma applies to the canonical trace
of the \emph{fixed} $w$. Choose a real compact smooth scalar
$\varphi$ equal to $|x|^2/2$ on a neighborhood of
$\overline B$. By properness $\varphi\circ t$ is an original
compact smooth scalar. Property (FP14) makes
$w-s(\varphi\circ t)$ actual psh for some $s>0$. Its canonical
Gauss trace is $g_w-s\varphi$ and is convex. Therefore
\[
                    \Delta g_w\ge2s\,dx\quad\text{on }B.
 \tag{FP17}
\]
Formula (FP4) supplies $g_w\in H_0^1(Q)$, and (FP5) identifies
$\tau=\Delta g_w=t_*\mu_w$. Apply the real lemma to get $k_j$
and $\epsilon_j$.

If some $u\in\mathscr C$ had source $f_j\mu_w$, its source
equality on the original tests $\varphi\circ t$ would imply
\[
 \Delta g_u=(1+\epsilon_j k_j)\Delta g_w,
                   \qquad g_u\in H_0^1(Q).
 \tag{FP18}
\]
Uniqueness of the real zero-boundary Poisson solution identifies
$g_u$ with the nonconvex $v_j$ in (FP9), contrary to the full
completed Gauss convexity theorem NF. This excludes every source
in the unchanged completion, not only toric or actual candidates.
The argument uses a subset of the full original test family to
derive a contradiction; it does not replace that family by a
smaller one in any definition.

Fixed-source domination FD gives all literal and strong assertions,
and its positive norm bound applied to $|f_j-1|\mu_w$ gives
(FP16). The bounds in (FP15) give equality of null ideals.
If $m=\|w\|_\infty>0$, then $1+w/m$ is an actual full normalized
tester and $\mu_w=m\mu_{1+w/m}$ in degree one. Hence
\[
       (f_j\mu_w)(A)\le\tfrac32 m C_S(A;V)
                  \le\tfrac32 m C_S^*(A;V)
                 \quad(A\text{ original Borel}).
 \tag{FP19}
\]
This uses no raw/outer equality. The same fixed measure bound
provides the common strong $L^1$ modulus from FD.

For the last assertion, choose balls $B_j\Subset Q$ escaping
every compact subset of $Q$, with closures still inside $Q$.
At each $B_j$ repeat (FP17) and the real lemma, choosing one
of its perturbations with amplitude below $1/j$. Each
$f_j-1$ has compact support
$t^{-1}(\operatorname{supp}k_j)\Subset V$, contained in
$t^{-1}(B_j)$ by properness.
Every compact analytic $H\subset V$ has compact image $t(H)$
in $Q$, disjoint from these supports for all sufficiently large
$j$. Thus $f_j-1$ is identically zero on an ordinary neighborhood
of $H$ eventually, as an original smooth germ after every actual
refinement; the zero representative is available there.
This gives convergence in the original smooth compact-open topology.
It does not claim that every ambient presentation of a zero germ
vanishes away from its actual analytic image.
Uniform convergence follows separately from the chosen amplitudes.
\end{proof}

\begin{corollary}[Degree-one density openness is already linear saturation]
For $q=1$ at any nonzero completed $w$ in the general original
scope of LI, existence of a two-sided $L^\infty(\mu_w)$ density
neighborhood of sources is equivalent to the known-source linear
saturation condition LI13. In particular LI13 fails at every fixed
actual $w$ of (FP14) in (FP1).
\end{corollary}

\noindent\textbf{Status: \texttt{DIRECT\_COROLLARY}} of LI,
degree-one source linearity and injectivity. This is part of the
same linearization mechanism, not another example family.

\begin{proof}
The reverse implication is LI with zero nonlinear remainder.
For the forward implication, shrink a density radius if necessary
and choose $0<a<1$ such that $(1\pm ag)\mu_w$ have sources
$v_\pm$ for every real $g$ with $\|g\|_\infty\le1$.
Degree-one additivity and source injectivity give
$v_++v_-=2w$. Thus
\[
 h=(v_+-w)/a,\qquad w\pm ah=v_\pm\in\mathscr C,
 \qquad M(h)=g\mu_w,
 \qquad\|h\|_w\le1/a.
 \tag{FP20}
\]
Homogeneity gives the bound for every bounded $g$. Linearized
injectivity in LI makes this solution a bounded linear inverse.
The last assertion follows from (FP15). No order-Lipschitz
assumption on a pre-existing nonlinear solution map is needed
for this degree-one argument.
\end{proof}

\subsection{Consequences, topology limits, and failed shortcuts}

The universally proposed implication
``a fixed reference admits both signs of all original compact
smooth tests, therefore every sufficiently small smooth density
change of its source is compatible'' is \texttt{FALSE} when
smallness means uniform density norm, even with support in one
fixed compact moment core. All data in this counterexample are
already literal strong data. It neither refutes LG nor uses a
nonliteral $L_\infty$ functional.

For the common-support construction, derivatives of the $k_j$ may
grow as their inner annuli shrink. It does \emph{not} assert
convergence in all smooth seminorms on that same compact core.
The smooth compact-open convergence in the last part instead uses
supports escaping the boundary and is not convergence in the
usual common-support topology of compact test functions.
Treating either of these stronger unsupported topologies as proved
would be a \texttt{FAILED\_ATTEMPT}.

The actual reference and domain never change in either sequence.
This differs from the already proved IS transfer, whose strict
references vary with the density parameter. For $q=1$, the
principal-ideal linear inverse hypothesis LI13 necessarily fails
at every $w$ in (FP14), since LI would otherwise solve all small
uniform density perturbations. This diagnoses a real geometric
obstruction to that sufficient condition, rather than silently
assuming a bounded-Hessian estimate from scalar ellipticity.

In fact the perturbations can be chosen entirely upward:
$\mu_w\le f_j\mu_w\le(1+\epsilon_j)\mu_w$.
Thus even this arbitrarily narrow one-sided Radon source bracket
at the same full-test-strict reference does not force compatibility.

The real $H_0^1$ boundary identification is essential. Distributional
equality of Laplacians alone would leave a harmonic summand and
would not prove the counterexample. Likewise the positivity of
$Kk_j$ is essential: it prevents an unexamined singular part of
the reference Laplacian from cancelling the logarithmic lower bound.
Both discarded shortcuts are recorded as \texttt{FAILED\_ATTEMPT}.

No general GQ/CTO statement, KC12 converse, general literal rooftop
theorem, or new capacity recovery conclusion follows. The independent
geometric Dirichlet compatibility problem remains \texttt{OPEN}.
All positive general implications elsewhere continue to use arbitrary
original middle-positive $S$ with their exact printed hypotheses.
No ordinary current identity here is promoted to all GK delta tests.

\paragraph{Exact input and external-source boundary.}
The internal sources are DO6--DO9 and SO5--SO6 for the prior
full-cone moment-Laplacian obstruction; NF6--NF11 for canonical
full completed convex traces;
GE7 and ED12 for positive quadratic energy and its completion;
HC14--HC15 and HD14 for original compact scalar decomposition and
the unnormalized full torus calibration; NP1--NP5 for the nonempty
all-test strict reference class; and FD1--FD8 for full literal
fixed-source domination. Their exact files, hashes, and read scopes
are in the companion ledger. The actual boundary-to-$H_0^1$ bridge,
its application to the full completed cone, and the fixed-source
obstruction above are new proofs in this round.

The one-dimensional almost-everywhere second expansion follows
from differentiation of the monotone derivative and the integral
formula for a convex function. A primary exposition of this
elementary one-dimensional input, and of the stronger theorem
which is unnecessary here, is Ralph Howard,
\emph{Alexandrov's theorem on the second derivatives of convex
functions via Rademacher's theorem}, May 1998, Remark 6.2,
printed page 19; the convex integral formula is recalled in Section 2.
The ordinary Green decomposition, weak Poisson uniqueness, maximum
principle and interior harmonic estimates are the standard real
facts documented in Jeff Viaclovsky's MIT 18.156 lecture notes,
Lectures 1--3 (interior estimates, Green decomposition and weak
harmonic regularity), together with Lecture 17 for the weak
Dirichlet problem via the Hilbert-space representation theorem.
Their basic real statements are not claimed as new.
The particular nonnegative angular-annular multiplier, its protection against
singular reference mass, and the original-category conclusions are
proved explicitly above. A targeted search of the exact frozen
local sources found no prior fixed all-test-strict-source version;
this is an internal novelty boundary, not a claim of world priority.
